\section{Problem Formulation}

\label{sec:problem}
% ══════════════════════════════════════════════════════════════════════════════

\subsection{Task Definition}

We consider the standard $N$-way $K$-shot classification problem. Given a
\emph{support set} $\Support = \{(x_i, y_i)\}_{i=1}^{N \times K}$
containing $K$ labeled images for each of $N$ novel species, and a
\emph{query set} $\Query = \{x_j\}_{j=1}^{Q}$ of unlabeled images, the
goal is to classify each query into one of the $N$ species.

\subsection{Taxonomic Structure}

Each species $s$ belongs to a taxonomic hierarchy:
\begin{equation}
  \text{Kingdom} \supset \text{Phylum} \supset \text{Class} \supset
  \text{Order} \supset \text{Family} \supset \text{Genus} \supset
  \text{Species}
\end{equation}

We denote the taxonomic label of species $s$ at rank $r$ as $t_r(s)$.
Two species in the same genus share $t_r$ for $r \in
\{\text{kingdom}, \ldots, \text{genus}\}$ but differ at species rank.

\subsection{Incorporating Priors}

For practical deployment, geographic coordinates $(\text{lat},
\text{lon})$ and observation date $d$ provide strong priors on which
species are plausible at a given location and time. We define the
candidate set:
\begin{equation}
  \mathcal{C}(\text{lat}, \text{lon}, d) = \{s : s \text{ has recorded
    occurrence within } r \text{ km and } \pm \Delta d \text{ days}\}
\end{equation}
which restricts the $N$-way classification to geographically and
temporally plausible species.


% ══════════════════════════════════════════════════════════════════════════════
